% !TEX root = ../5.DISTILLATION.tex
\section{DISCUSSION}

\subsection{Effect of reflux ratio on product purity and column performance}

Based on experiments 1, 2, and 3, it is evident that as the reflux ratio ($R$) increases, the purity of the product (distillate) also increases. This phenomenon can be explained by the energy balance in the column. With a constant heat supply, the total molar vapor flow ($V$) remains relatively constant.

Since $V = L + D$, an increase in the reflux flow ($L$) results in a decrease in the distillate product flow ($D$). For the distillation column, the heat supplied to the reboiler serves to separate the phases. When the product flow $D$ decreases, the specific heat supplied per mole of product increases, leading to better separation and higher purity in both top and bottom products.

However, it must be noted that general performance does not allow for a complete assessment of the economic viability of the equipment. While an increased reflux flow improves the separation efficiency of each tray—potentially reducing the total number of trays required and lowering initial capital costs—it also increases the individual heat supply requirement per unit of product. This energy cost can account for a significant percentage of the total operating costs, and must be balanced against the desired product purity.

\subsection{Effect of feed location on purity and performance}

The location of the feed tray significantly influences the concentration profile within the column. As the vapor and liquid streams pass through each tray, their concentrations change based on the exchange capacity of the tray.

Inputting the feed at trays near the bottom increases the number of rectification stages above the feed, generally increasing the distillate purity. Conversely, feeding near the top reduces the rectification stages. However, if the physical feed position deviates significantly from the optimal theoretical feed tray, the overall efficiency of the trays decreases. Experimental results show that the input tray position does not directly affect the individual tray efficiency, but rather how closely the column operates to its theoretical potential.

\subsection{Stable operating conditions}

Stable operation is achieved when the heat exchangers ensure consistent heat supply and the flow meters show steady values. Each tray maintains a distinct boiling temperature, which decreases progressively from the bottom (reboiler) to the top (condenser). To maintain stability, the following are required:
\begin{itemize}
    \item Proper operation of the heat exchangers.
    \item Stable flow rates for feed, reflux, and distillate.
    \item Consistent temperature measurements across the trays (higher trays indicate lower temperatures).
\end{itemize}

\subsection{Error analysis and mitigation}

Sources of error during the experiment include:
\begin{itemize}
    \item Unstable flow rates due to valves not being fully open or mechanical fluctuations.
    \item \textbf{Rotameter instability:} The observation that "marbles" (floats) are dropped or oscillate, requiring constant monitoring of the flow.
    \item Measurement inaccuracies in alcohol concentration using hydrometers ("alcohol by incorrect design").
    \item Timing discrepancies when reading multiple sensors simultaneously.
    \item Data rounding and plotting inaccuracies during the calculation phase.
\end{itemize}

To mitigate these errors, measurements should be taken simultaneously, flow rates must be monitored and adjusted continuously, and calculations should be verified using standardized tables and software where applicable.
